\section{Amplification against fast travel}\label{sec:fast}
The fixed-height renewal theorem gives one polynomial failure power for a
short bridge. We need every polynomial power in order to sum over the
positions and outside pieces of arbitrary subpaths. The following argument
uses a large but \emph{fixed} number of turns, chosen after the desired
failure exponent. Adjacent-arm avoidance pays for most turns; a long
irreducible shared by the two ends is integrated once.

\subsection{A length deficit and recoverable turns}

For an event \(E\) of height-\(h\) bridges, write \(B_h[E]\) for
its unnormalized critical mass.  A weak bridge is a path whose endpoints
attain its extreme heights; intermediate contacts with their levels
are allowed.

\begin{lemma}[An integrated short-bridge estimate]\label{09_fast_paths:integrated-short}
For every sufficiently small \(\varepsilon>0\), there are
\(\kappa,a_0>0\) such that, for each fixed \(C<\infty\),
\begin{equation}
 \sum_{0\le j\le Cr}B_j[\,|\gamma|\le2H^{4/3-\varepsilon}\,]
       \le C' r^{3/4}H^{-a_0},
 \qquad H^{1-\kappa}\le r\le CH.                 \label{eq:09-short-gain}
\end{equation}
\end{lemma}
\begin{proof}
The polynomial lower-length failure bound in
Theorem~\ref{thm:08-bridge-length} gives, for some \(a>0\),
\(B_j[|\gamma|\le j^{4/3-\varepsilon/2}]
\le C B_jj^{-a}\).  Choose \(\kappa>0\) so small that
\((1-2\kappa)(4/3-\varepsilon/2)>4/3-\varepsilon\).
For large \(H\), if \(j\ge H^{1-2\kappa}\), the length cap in
Equation~\eqref{eq:09-short-gain} is below that threshold.  Such
heights contribute at most
\(Cr^{3/4}H^{-a(1-2\kappa)}\).  Smaller heights, including zero,
have total mass \(O(H^{3(1-2\kappa)/4})\), at most
\(Cr^{3/4}H^{-3\kappa/4}\) in the stated range.  Taking
\(a_0\le\min\{a(1-2\kappa),3\kappa/4\}\) proves the lemma.
\end{proof}

\begin{lemma}[Turning-extremum decomposition]\label{09_fast_paths:turns}
Fix \(d\ge1\) and put \(k=2^d\).  A weak up bridge of height in
\([H,2H]\) admits an averaged decomposition into \(k\) weak up
pieces, each of span in \([c_dH,2H]\), separated by \(k-1\) retreats.
If the retreat spans have dyadic scales \(s_\ell\), with scale 1
including empty and bounded retreats, the total level-choice weight
of a fixed decomposition is at most
\(C_d\prod_\ell(s_\ell/H)\).

Every piece can be recoded as a strict port bridge, with bounded weight
change and finitely many local inverse records per cut.  It gains only
\(O(1)\) length.  At a turn whose retreat span grows as a positive
power of \(H\), the two departing arms start at neighboring ports
and remain disjoint through first visits to any common distance
\(o(s_\ell)\) and \(o(H)\) into their shared side.
\end{lemma}
\begin{proof}
Average a cut level \(m\) over lattice lines in the middle third of
the bridge.  If the line is crossed once, cut there and insert an empty
retreat.  Otherwise, after the first hit of \(m\), let \(i\) be the
last point attaining the minimum height of the remaining path, and
let \(j\) be a point attaining the maximum height before \(i\),
with a fixed tie convention.  The multiple crossing implies
\(i<m<j\), and \(j\) occurs after the first hit.  The three pieces
are the up bridge to \(j\), the down retreat to \(i\), and the
remaining up bridge.  They are weak bridges: the defining extrema
bound all their intermediate heights.  Repeat on both up pieces
through depth \(d\).  Each up span loses at most a fixed factor per
generation, so its final span is at least \(c_dH\).

For a fixed final decomposition, the cut at an internal tree node
must lie in its retreat's height interval, or at the single-crossing
cut if the retreat is empty.  Its number of choices is \(O(s_\ell)\).
The averaging denominator at that node is \(\asymp H\), with
constants depending on \(d\).  The tree order is fixed, so multiplying
these bounds proves the asserted placement factor.

We describe the local recoding because neighboring starts are needed
for Corollary~\ref{09_fast_paths:adjacent-arms}.  A nonempty peak is a
down-pointing triangle, two thirds up its band, visited via its two
lower neighbors.  Assign that triangle to the retreat, entering through
its top port.  The reversed up piece can start at the neighboring top
port on its branch side and pass through the adjacent down-pointing
triangle to the same lower neighbor of the peak.  If it already used
that triangle, the visit must be immediately next to this lower
neighbor: it cannot cross above its extremum, and the piece continues
macroscopically below.  In that case start directly there and remove
the short peak portion.  The adjacent triangle cannot be in the
relevant initial retreat arc, because any such visit would use that
same lower neighbor, contradicting the original disjointness.

Reflect this construction at minima.  If an original endpoint is a
vertex, extend it to the nearest exterior band line by the bounded
construction in the proof of Lemma~\ref{lem:port-extension}.
The added vertices lie outside the old height range; the extension
changes length and height by \(O(1)\), has bounded multiplicity, and
changes weight by a bounded factor.  Original port endpoints and
single-crossing port cuts already lie on boundary lines.  These local
operations produce strict port bridges within the outer band lines of
their extrema.  At a macroscopic turn, modifications at remote ends
are outside the first-hit portions of distance \(o(s_\ell)\) and
\(o(H)\); hence those neighboring arms remain disjoint.  A minimum
attained again at the starting height is still a local minimum between
upper neighbors, so the same construction applies there.  Each cut
has only finitely many possible local records and weight ratios.
Relative shapes and these records recover the original pieces, their
translations being fixed by concatenation up to bounded lattice
choices.  In upper bounds all ending heights may therefore be summed
freely over their indicated scale ranges.
\end{proof}

\subsection{Integrating stopped and terminally marked ends}
The local path modification has made the two arms at a turn start from
neighboring ports. Their endpoint laws still differ: one arm may stop at
first passage, while the other is summed over a terminal renewal. The next
lemma isolates that change of measure before any avoidance factors are
multiplied.

\begin{lemma}[Avoidance with marked terminal renewals]\label{09_fast_paths:marked-conditioning}
Consider two independent renewal end strings from neighboring ports.
For either string use one of the following measures: its probability
law stopped at the first renewal above \(q\), or its bridge measure
summed over a marked terminal renewal with
\(f\le1+U<2f\), where \(f\ge1\).  Call the first kind good and the
second kind marked.  A good end has mass scale 1 and probe scale \(q\);
a marked end has mass scale \(f^{3/4}\) and probe scale \(f\).
Requiring the two strings to avoid through their first visits to a
small fixed multiple of the smaller probe scale costs at most
\(S_2(t)\), times a constant and the two mass scales.  Here \(t\)
is that small multiple, rounded to an integer; bounded scales use
the bound 1.  The assertion remains valid at all turns of a chain
simultaneously when different turns use distinct end strings and all
other incidence constraints are dropped.
\end{lemma}
\begin{proof}
For unbounded scales choose \(t\) less than, say, one tenth of both
probe scales, reducing the constant to handle fixed scale conventions.
Expose each infinite end walk through its first renewal at or above
\(t\).  The first-visit avoidance event is measurable with respect
to these two stopped strings.  A good string has conditional remaining
mass at most 1.  For a marked string, condition on its exposed endpoint
height \(a\).  If \(a\ge2f\) there is no admissible marker, except
for harmless bounded boundary cases.  Otherwise its conditional
expected number of possible marked renewals is at most
\[
                       1+\sum_{j\le2f}B_j\le Cf^{3/4}.
\]
This includes the stopping renewal itself and is uniform in overshoot.
The infinite renewal law up to a marked renewal has exactly the
unnormalized bridge weights, so the argument bounds the required
finite-string measure.  After multiplying the two conditional bounds,
the remaining event has probability \(S_2(t)\).  This proves all
good/marked combinations without conditioning a terminal height in
advance.

For simultaneous turns, first express every end in relative coordinates.
The local records at a turn fix its neighboring starting ports up to a
bounded choice.  Once middle bridges and large jumps have been summed,
discard consistency of their translations and heights with these
relative strings.  This enlarges the set being counted.  Each end then
appears in exactly one turn, so the pairs of independent string variables
used by different avoidance tests are disjoint.  The preceding
conditional integration can be performed at every turn, and its
factors multiply.  In particular, no avoidance estimate is being
applied to a string already biased by a different turn.
\end{proof}

\subsection{The gain from many turns}
We now combine the short-bridge deficit with adjacent-arm avoidance.
The partition into disjoint atomic end strings is part of the proof: it
is what permits multiplication of turn costs without conditioning one
test on the success of another.

\begin{theorem}[Fast bridges]\label{09_fast_paths:fast-bridges}
For every \(\varepsilon>0\) and every \(A<\infty\),
\begin{equation}
 \sum_{H\le h\le2H}
 B_h[\,|\gamma|\le H^{4/3-\varepsilon}\,]
       \le C_{\varepsilon,A}H^{-A}.
                  \label{eq:F-33b}
\end{equation}
\end{theorem}
\begin{proof}
It suffices to take small \(\varepsilon>0\).  Use
Lemma~\ref{09_fast_paths:turns} with \(k=2^d\), where \(d\) is
fixed now and chosen large at the end.  Local recodings increase the
total length by \(O(k)\), so every resulting portion has length at
most \(2H^{4/3-\varepsilon}\) for large \(H\).  Let
\(\kappa,a_0\) come from Lemma~\ref{09_fast_paths:integrated-short},
and choose
\[
           0<\Delta<\min\{1/40,\kappa/10,a_0/10\}.
\]
An up piece has scale \(r=H\), while a retreat has scale
\(r=s_\ell\).  The unrestricted summed-height mass of either is
\(O(r^{3/4})\).  Retreats with \(s_\ell\le H^{1/20}\) will use
only this bound.

For any other piece inspect from both ends its first renewal beyond
\(q=\lfloor rH^{-\Delta}\rfloor\), a threshold less than half
its span.  The irreducible decomposition is unchanged by reversal.
If the two stopped strings do not overlap, allowing a shared cut
level, call the piece type \(\mathrm M\).  Sum its two end strings
separately and give its middle bridge the full mass bound
\(Cr^{3/4}\).  If \(r\ge H^{1-\kappa}\), the middle also has
the length cap, so Equation~\eqref{eq:09-short-gain} supplies the
extra factor \(H^{-a_0}\), regardless of its actual remaining span.

An \(\mathrm M\) end is good, denoted \(\mathrm G\), if its
stopping jump lands within distance \(cr\) of that end, for a small
fixed \(c>0\).  Ignore this extra constraint and use its stopped
renewal probability law.  Otherwise the end is bad.  Group
\(1+\)the height before its stopping jump in a dyadic bin of scale
\(f\le Cq\).  Its preceding marked-renewal prefix has mass
\(O(f^{3/4})\), and its next independent jump has height at least
\(cr/2\), with probability \(O(r^{-3/4})\).  Thus the bad end
costs \(C(f/r)^{3/4}\) in addition to the middle's baseline.

If the end strings overlap, call the piece type \(\mathrm Z\).
They contain the same irreducible jump between the two near-end
ranges.  There cannot be two such jumps, since the ranges are
disjoint and the irreducibles are linearly ordered in height.  Mark
both ends bad.  Summing their preceding prefixes and this single
shared jump gives
\begin{equation}
 C f_1^{3/4}f_2^{3/4}r^{-3/4}
       =Cr^{3/4}(f_1/r)^{3/4}(f_2/r)^{3/4}.        \label{eq:09-overlap}
\end{equation}
The shared jump has height comparable to \(r\).  There is no
length gain for this piece, and its tail is charged only once.

For each turn of a retreat with \(s>H^{1/20}\), impose adjacent-arm
avoidance at a small multiple of the minimum probe scale: \(q\) at
a good end, \(f\) at a bad end.  These distances are \(o(s)\) and
\(o(H)\), as required by the local recoding.  When the minimum is
bounded use the constant bound instead.  By
Lemma~\ref{09_fast_paths:marked-conditioning}, the marked terminal
renewal sum after exposing this avoidance still costs at most
\(Cf^{3/4}\).  Large-jump tail factors follow these prefix portions
and do not consume them.  In the \(\mathrm Z\) case the one shared
jump has already been integrated as in Equation~\eqref{eq:09-overlap}.
The remaining end variables are distinct at distinct turns, so their
avoidance gains multiply.  Corollary~\ref{09_fast_paths:adjacent-arms}
therefore supplies the minimum probe scale to power
\(-3/4+o(1)\) at every such turn.

Each turn adjoins an up piece of scale \(H\) to a retreat of scale
\(s\).  Combine its avoidance factor with its attached bad-end
costs.  Apart from \(H^{o(1)}\), it supplies \(s^{-3/4}\) and
the following additional factors:
\[
\begin{array}{c|c}
 \text{end types}&\text{factor after extracting }s^{-3/4}\\ \hline
 \mathrm G,\mathrm G&H^{3\Delta/4}\\
 \text{one bad}&1\\
 \text{both bad}&H^{-3\Delta/4}.
\end{array}
\]
For two good ends, the smaller threshold is comparable to
\(sH^{-\Delta}\).  With one bad end of scale \(r\), if its probe
\(f\) is smaller, its \(f^{3/4}\) cost cancels the avoidance
factor, leaving \(r^{-3/4}\le Cs^{-3/4}\).  If the good probe
\(r'H^{-\Delta}\) is smaller, use \(f/r\le CH^{-\Delta}\)
and \(s\le Cr'\).  With two bad ends, writing their probes as
\(f_1,f_2\), the factor is exactly
\[
 (f_1/H)^{3/4}(f_2/s)^{3/4}
      \min(f_1,f_2)^{-3/4}
   =s^{-3/4}\bigl(\max(f_1,f_2)/H\bigr)^{3/4},
\]
which is at most \(Cs^{-3/4}H^{-3\Delta/4}\).

Nominally extract \(s_\ell^{-3/2}\) from the two turns of every
retreat.  Multiplying all baseline masses and placement factors gives
\begin{equation}
 H^{3k/4}\prod_{\ell=1}^{k-1}
  \left[s_\ell^{3/4}\frac{s_\ell}{H}s_\ell^{-3/2}\right]
           =H^{3/4}\prod_{\ell=1}^{k-1}(s_\ell/H)^{1/4}.
                                                        \label{eq:09-chain}
\end{equation}
For an untreated retreat \(s\le H^{1/20}\), restore the nominal
factor \(s^{3/2}\).  Its net factor is at most
\(H^{-1/4}s^{7/4}\le H^{-13/80}\).  A treated retreat with
\(s<H^{1-\kappa}\) pays at least \(H^{-\kappa/4}\) through
its scale factor in Equation~\eqref{eq:09-chain}.

Every both-good loss can be charged to its incident \(\mathrm M\)
up piece, since a \(\mathrm Z\) piece has both ends bad.  Each up
piece has at most two turns, so its own \(H^{-a_0}\) gain leaves
at least \(H^{-a_0/2}\), by the choice of \(\Delta\).  Assign
each \(\mathrm Z\) up piece to one incident retreat; an endpoint
up piece still has one such retreat because \(k\ge2\).  A retreat
receives at most two assignments.  If its scale is below
\(H^{1-\kappa}\), its scale saving pays those assignments.  If
it is larger and type \(\mathrm M\), its own \(H^{-a_0}\) length
gain pays them.  If it is larger and type \(\mathrm Z\), the
assigned turn is both bad and pays \(H^{-3\Delta/4}\).  Thus there
is a constant \(b=b(\varepsilon)>0\), independent of fixed \(k\),
such that all up pieces together save \(H^{-bk}\).  For example
one can choose any sufficiently small positive number below
\(\min\{a_0/2,\kappa/8,13/160,3\Delta/8\}\).

For fixed \(k\), there are finitely many type choices and only
logarithmically many dyadic scales per piece.  All avoidance losses
are \(H^{o(1)}\) for this fixed number of tests.  The total mass
is consequently at most \(C_kH^{3/4-bk+o(1)}\).  Taking a fixed
power of two \(k\) sufficiently large in terms of \(A\), and
then enlarging the constant for bounded \(H\), proves
Equation~\eqref{eq:F-33b}.
\end{proof}

